\chapter{ไฟฟ้าสถิต สนามไฟฟ้า และศักย์ไฟฟ้า}
\label{ch:ch11_electrostatics}
\index{ไฟฟ้าสถิตและสนามไฟฟ้า}

\begin{conceptbox}[title=วัตถุประสงค์การเรียนรู้ประจำบท]
เมื่อสิ้นสุดการศึกษาบทเรียนนี้ นักศึกษาสามารถ
\begin{enumerate}
    \item เมื่อศึกษาบทนี้จบแล้ว ผู้เรียนควรสามารถ
    \item อธิบายธรรมชาติของประจุไฟฟ้า ชนิดของประจุ และหลักการอนุรักษ์ประจุไฟฟ้าได้
    \item อธิบายกระบวนการทำให้วัตถุเกิดประจุไฟฟ้าด้วยการเสียดสี การสัมผัส และการเหนี่ยวนำ
    \item คำนวณแรงระหว่างประจุไฟฟ้าโดยใช้กฎของคูลอมบ์ได้อย่างถูกต้อง
    \item คำนวณสนามไฟฟ้าและศักย์ไฟฟ้าที่เกิดจากประจุจุดในสถานการณ์ต่าง ๆ ได้
\end{enumerate}
\end{conceptbox}

\section{ทฤษฎีและหลักการพื้นฐาน}

\begin{bioappbox}
\textbf{ชีวฟิสิกส์ของศักย์เยื่อหุ้มเซลล์และสมการเนิร์นสต์}

เยื่อหุ้มเซลล์ของสิ่งมีชีวิตทำหน้าที่เป็นตัวเก็บประจุชีวภาพ (Biological Capacitor) โดยมีสารฟอสโฟลิพิดสองชั้นคั่นระหว่างไอออนบวกและไอออนลบภายในและภายนอกเซลล์ ทำให้เกิดความต่างศักย์เยื่อเซลล์ขณะพัก (Resting Membrane Potential) ประมาณ $-70\,\text{mV}$ ซึ่งสามารถคำนวณได้อย่างแม่นยำด้วยสมการเนิร์นสต์ (Nernst Equation) $V = \frac{RT}{zF} \ln\frac{[C_{\text{out}}]}{[C_{\text{in}}]}$ ความรู้นี้เป็นรากฐานของสรีรวิทยาระบบประสาทและกล้ามเนื้อ
\end{bioappbox}

\begin{labbox}
1. ตรวจสอบตำแหน่งศูนย์ (Zero Error) ของเครื่องมือวัดทุกชนิดก่อนเริ่มบันทึกข้อมูล
2. ทำการวัดซ้ำอย่างน้อย 3 ถึง 5 ครั้งในแต่ละจุดทดลองเพื่อคำนวณหาค่าเฉลี่ยและค่าเบี่ยงเบนมาตรฐาน
3. บันทึกผลการวัดพร้อมระบุหน่วยเอสไอ (SI Units) และค่าความคลาดเคลื่อนของการวัดเสมอ
\end{labbox}

\section{ตัวอย่างโจทย์และการวิเคราะห์เชิงฟิสิกส์}

\begin{examplebox}
ตัวอย่างที่ 11.1 ประจุไฟฟ้าสองประจุ q₁ = 2.0$\times$10$^-$⁶ C และ q₂ = 3.0$\times$10$^-$⁶ C วางห่างกัน 0.30 m จงหาขนาดของแรงระหว่างประจุทั้งสอง

วิธีทำ ใช้กฎของคูลอมบ์

F = kq₁q₂/r$^2$

F = (8.99$\times$10⁹)(2.0$\times$10$^-$⁶)(3.0$\times$10$^-$⁶)/(0.30)$^2$
\end{examplebox}

\section{แบบฝึกหัดทบทวนและคำถามท้ายบท}

\begin{enumerate}
    \item จงอธิบายความแตกต่างระหว่างการทำให้วัตถุเกิดประจุด้วยการสัมผัสและด้วยการเหนี่ยวนำ
    \item ประจุไฟฟ้าสองประจุขนาด 4.0$\times$10$^-$⁶ C และ 6.0$\times$10$^-$⁶ C วางห่างกัน 0.20 m จงหาขนาดของแรงระหว่างประจุทั้งสอง
    \item หากระยะห่างระหว่างประจุสองประจุเพิ่มขึ้นเป็นสามเท่าจากเดิม แรงระหว่างประจุจะเปลี่ยนแปลงไปอย่างไร
    \item จงหาสนามไฟฟ้า ณ ตำแหน่งที่ห่างจากประจุจุด 8.0$\times$10$^-$⁹ C เป็นระยะ 0.10 m
    \item สนามไฟฟ้าที่ตำแหน่งหนึ่งมีขนาด 500 N/C จงหาแรงที่กระทำต่อประจุทดสอบขนาด 2.0$\times$10$^-$⁶ C ที่วางอยู่ ณ ตำแหน่งนั้น
    \item จงอธิบายความหมายของเส้นแรงไฟฟ้า พร้อมระบุคุณสมบัติสำคัญสามประการของเส้นแรงไฟฟ้า
\end{enumerate}

% สิ้นสุดบทเรียน
